\documentclass[11pt]{article}
\usepackage[margin=1in]{geometry}
\usepackage{amsmath,amssymb,amsthm,mathtools}
\usepackage[hidelinks]{hyperref}
\newtheorem{theorem}{Theorem}
\newtheorem{lemma}[theorem]{Lemma}
\newtheorem{corollary}[theorem]{Corollary}
\title{All-index cubic Strong Factorial by adjoint half-plane separation}
\author{Internal proof and exact replay; no priority claim}
\date{17 November 2026}
\begin{document}
\maketitle
For $(\gamma)_j=\gamma(\gamma+1)\cdots(\gamma+j-1)$ put
\[
 T_{n,\gamma}(t)=\sum_{k=0}^n\binom nk(\gamma)_{2k}t^k.
\]
\begin{theorem}\label{thm:simple}
If $n\ge1$, $\gamma>0$ and $\gamma\ge n-2$, then $T_{n,\gamma}$
has only simple complex roots.
\end{theorem}
\begin{corollary}\label{cor:sfc}
For every positive Gamma shape $\beta$, every nonzero
$f(X)=X(A+BX^2)$ with complex coefficients has no two consecutive
zero moments. This holds at every starting index. Multiplication by
monomials in other independent variables preserves the assertion,
with arbitrary positive shapes in those variables.
\end{corollary}
The proof is uniform: there is no asymptotic cutoff or finite-prefix
obligation. It closes this cubic family, not unrestricted Strong Factorial
or the general Factorial Conjecture. The previously retained prime-power
index and fold results are independent earlier-stage evidence.

\section{A half-plane estimate and its convolution consequence}
Define the monic Gamma moment polynomial
\[
 G_{n,\alpha}(z)=\sum_{k=0}^n\binom nk(\alpha)_kz^{n-k},\qquad \alpha>0.
\]
\begin{lemma}\label{lem:gamma}
Every zero of $G_{n,\alpha}$ has real part at most
$-(\alpha-n+1)$. The inequality is strict when $n\ge2$.
\end{lemma}
\begin{proof}
Put $S=\alpha+n-1$ and form a complex symmetric tridiagonal matrix
$J$, with entries $J_{jj}=S+2-2j$ and
$J_{j,j+1}=J_{j+1,j}=i\sqrt{j(S+1-j)}$, $1\le j<n$.
The radicands are positive. Its leading principal characteristic
polynomials obey
\[
 D_0=1,\quad D_1=z-S,\quad
 D_j=(z-S-2+2j)D_{j-1}+(j-1)(S+2-j)D_{j-2}.
\]
The coefficient formula
\[
 D_j(z)=\sum_{k=0}^j(-1)^{j-k}\frac{j!}{k!}
                         \binom{S-k}{j-k}z^k
\]
verifies this recurrence directly. At $j=n$ it gives
$G_{n,\alpha}(z)=\det(zI+J)$, also for noninteger $\alpha$.
The Hermitian part of $J$ is its real diagonal, whose smallest entry
is $\alpha-n+1$. Thus each eigenvalue of $-J$ has real part at most
the negative of this entry. Equality for $n\ge2$ would require an
eigenvector supported solely in its last coordinate. The nonzero last
off-diagonal entry rules that out.
\end{proof}

We use the following classical input, in precisely its half-plane case:
the Grace--Walsh--Szeg\H{o} coincidence theorem says that the symmetric
multiaffine polarization of a degree-$n$ polynomial nonvanishing on an
open half-plane is nonvanishing whenever all $n$ arguments belong to
that half-plane. A source with statement and proof is Wagner,
\emph{Multivariate stable polynomials: theory and applications},
Theorem 4.1 and Section 4, in the
\href{sources/wagner-multivariate-stable-polynomials-2011.pdf}{cached author PDF}
(\href{https://www.math.uwaterloo.ca/~dgwagner/ceb_wagner.pdf}{author URL}).
Only this stated classical input is imported; the application below is explicit.

For monic degree-$n$ polynomials define their additive convolution by
\[
 (f\boxplus g)(z)=\frac1{n!}\sum_{k=0}^n f^{(k)}(z)g^{(n-k)}(0).
                                                               \tag{1}
\]
If $s_1,\ldots,s_n$ are the roots of $g$ and $F$ is the symmetric
multiaffine polarization of $f$, coefficient comparison gives
\[
 (f\boxplus g)(z)=F(z-s_1,\ldots,z-s_n).                         \tag{2}
\]
In particular, if $f$ has no roots in $\Re z>A$ and all roots of $g$
have $\Re s<B$, then $f\boxplus g$ has no roots in $\Re z\ge A+B$:
all the arguments in (2) then belong to $\Re z>A$, so the classical
input applies. This uses polarization, not closure of a convex average
of stable polynomials.

Let $U,V$ be independent unit-rate Gamma variables of shapes
$\gamma$ and $n+1$. Define
\[
 B(z)=\mathbb E(z+U+V/2)^n
     =\sum_{j=0}^n n_{\underline j}c_j z^{n-j},\qquad
 \sum_{j\ge0}c_jx^j=(1-x)^{-\gamma}(1-x/2)^{-(n+1)}.             \tag{3}
\]
Here $n_{\underline j}=n!/(n-j)!$. This is
$G_{n,\gamma}\boxplus [2^{-n}G_{n,n+1}(2z)]$, by expansion of
(1), or by independence of $U,V$. Lemma \ref{lem:gamma} and (2) show
\[
 \Re z<n-\gamma-2\quad\hbox{at every root of $B$, if $n\ge2$}.   \tag{4}
\]
For $n=1$, $B(z)=z+\gamma+1$ has a strictly negative root directly.
Under the hypotheses of Theorem \ref{thm:simple}, $B$ is therefore
strictly Hurwitz: every root has negative real part. Consequently
$B(z)$ and $B(-z)$ have disjoint zero sets.

\section{The exact differential and adjoint equations}
Put $M=\gamma+2n$, and reverse $T$ after a quadratic substitution:
\[
 R(z)=z^{2n}T_{n,\gamma}(-z^{-2}).                              \tag{5}
\]
It is monic with $R(0)=(-1)^n(\gamma)_{2n}\ne0$.
The coefficient recurrence of $T$ gives, with $\theta=tD_t$,
\[
 [D_t+(\theta-n)(2\theta+\gamma)(2\theta+\gamma+1)]T=0.
\]
Putting $\delta=zD_z$ in (5) transforms this into
\[
 [\delta(\delta-M)(\delta-M-1)-z^2(\delta-2n)]R=0,
\]
or $LR=0$, where
\[
 L=D_z^3+aD_z^2+bD_z+c,\quad
 a=\frac{2-2M}{z},\quad b=\frac{M(M-1)}{z^2}-1,\quad c=\frac{2n}{z}.
                                                               \tag{6}
\]
All nonzero finite points are ordinary points. A repeated root of $R$
at such a point has multiplicity exactly two: a multiplicity $\ge3$
would give an uncancelled lowest Taylor term in the monic third derivative.

The generating function $C(x)$ in (3) satisfies
\[
 (2-3x+x^2)C'(x)=[(2\gamma+n+1)-(\gamma+n+1)x]C(x).
\]
Equivalently, with $c_{-1}=0$,
\[
 2(j+1)c_{j+1}=(3j+2M-3n+1)c_j-(j+M-n)c_{j-1}.
\]
Substitution in (3), including both end coefficients, gives
\[
 \mathcal EB:=zB'''-(3z+M-1)B''+(2z+2M-2)B'-2nB=0.            \tag{7}
\]
The formal adjoint of (6) is
\[
 L^*f=-f'''+(af)''-(bf)'+cf.
\]
On any local branch away from zero, let
\[
 f_+(z)=e^{-z}z^{1-M}B(z),\qquad
 f_-(z)=e^{z}z^{1-M}B(-z).                                    \tag{8}
\]
Both solve $L^*f=0$. For $f_+$ the exact gauge identity is
\[
 L^*(e^{-z}z^{1-M}B)=-\frac{e^{-z}z^{1-M}}z\mathcal EB.
                                                               \tag{9}
\]
It follows by differentiating with logarithmic derivative
$-1+(1-M)/z$. Reflection gives $f_-$, because $a,c$ are odd and $b$
is even. Powers with noninteger $M$ need only be defined locally;
reflection changes the chosen power by a nonzero constant.

\section{The zero concomitant rules out every repeated root}
For a solution of $LR=0$ and a solution of $L^*f=0$, the Lagrange
concomitant
\[
 \mathcal C(f,R)=fR''+(af-f')R'+[f''-(af)'+bf]R                \tag{10}
\]
has derivative $fLR-RL^*f=0$. In (8), each value of (10) is
$e^{\mp z}z^{1-M}V_\pm(z)$ for a rational function $V_\pm$.
It must in fact vanish identically. Otherwise its constancy implies
\[
 V_\pm'/V_\pm=\pm1+(M-1)/z,
\]
an impossible rational identity: the left side tends to zero at infinity,
whereas the right side tends to $1$ or $-1$. This argument works equally
for integer and noninteger $\gamma$; it does not assert that $f_\pm$
itself is rational.

If $z_0$ were a repeated root of $R$, then $z_0\ne0$ and
$R(z_0)=R'(z_0)=0$, $R''(z_0)\ne0$. Evaluating the two zero
concomitants gives
\[
 B(z_0)=B(-z_0)=0,
\]
contradicting the half-plane separation (4). Thus $R$ is squarefree.
The substitution $t=-z^{-2}$ is unramified away from zero and
$T(0)=1$, so a repeated root of $T$ would be a repeated root of $R$.
This proves Theorem \ref{thm:simple}.

\section{Consecutive Gamma moments}
For terminal order $n\ge1$ put $\gamma=n+\beta-1>0$ and
$\vartheta=tD_t$. Exact coefficient identities give
\[
 T_{n,\gamma+1}=T_{n,\gamma}+\frac2\gamma\vartheta T_{n,\gamma},
 \qquad
 T_{n-1,\gamma}=T_{n,\gamma}-\frac1n\vartheta T_{n,\gamma}.      \tag{11}
\]
The determinant of this two-by-two linear system is
$-(\gamma+2n)/(\gamma n)\ne0$.
For $AB\ne0$, the moments at orders $n,n-1$ of $X(A+BX^2)$ are
nonzero scalar multiples of $T_{n,\gamma+1}(B/A)$ and
$T_{n-1,\gamma}(B/A)$, respectively. Since $\gamma>n-1$, Theorem
\ref{thm:simple} and (11) exclude simultaneous zeros. Monomial cases
are immediate, and independent-variable monomials have nonzero Gamma
moments at every order. This proves Corollary \ref{cor:sfc}.

\section{Verification and boundaries}
The paired module and checker replay the exact differential, gauge,
concomitant, convolution, matrix, and moment identities. Symbolic
operator identities check arbitrary parameters; bounded rational cases
are regression controls, not the proof of the universal claim.
The classical input is pinned by a source hash and its theorem-level
application was reviewed. Independent replay, formal verification,
external review and novelty claims are all false.
No statement for arbitrary $\gamma>0$ below $n-2$, all lower exponents,
all gaps, general supports or unrestricted FC/SFC follows here.
\end{document}
